% ---
% id: reductions/thm:factorization
% kind: thm
% title: "Equality of the physical and squashed output states"
% summary: >-
%   Replacing each round's physical measurement by the squashing channel
%   followed by the target measurement leaves the protocol's whole output state
%   unchanged. Every security criterion therefore takes the same value on the
%   physical and the squashed protocol, with no error term.
% prominence: standard
% status: accepted
% uses: [def:protocol, def:criterion, lem:msmfact]
% assumes: [A-MEAS, A-ROUND-SQ, A-SEQ]
% refs: []
% provenance: newly-derived
% lean:
% ---

A \tref{squashing relation} is a statement about one round's measurement.
Through the factorization of the \tref{measurement channel}
(Lemma~\ref{lem:msmfact}) it
becomes a statement about the whole protocol, and the proof uses no property
of the criterion $f_E$ beyond its being a function of the output state: not
linearity, not monotonicity, not continuity. The one structural hypothesis is
the factorization $\Grnd_j=\bar{\Grnd}_j\circ(\idmap\otimes\redNmsm_M)$ of
Definition~\ref{def:protocol}, and it rests on the measuring devices
entering only through one POVM per round (Assumption~\ref{ass:meas})
together with the round order and the one POVM for all rounds that the
definition declares. The round's order is not part of
Assumption~\ref{ass:meas}: that every public message of the round follows
every measurement of the round is A-SEQ (\xref{transmitter/ass:seq}).

\begin{theorem}
\label{thm:factorization}
Let $M\dsq_{\Lambda}F$ over $\Vset$ and let $\Eprot_M,\Eprot_F$ be as in
Definition~\ref{def:protocol}. Then
\begin{equation}\label{eq:factorization}
  \Eprot_M\;=\;\Eprot_F\circ\big(\idmap\otimes\Lambda^{\otimes n}\big),
\end{equation}
the identity acting on every register other than $\Bsys_1\dots\Bsys_n$; and
consequently, for every finite-dimensional adversary system $E$ and every input
$\rho$,
\begin{equation}\label{eq:sameoutput}
  (\Eprot_M\otimes\idmap_E)(\rho)
  =(\Eprot_F\otimes\idmap_E)\big(\sigma\big),
  \qquad
  \sigma\coloneqq\big(\idmap\otimes\Lambda^{\otimes n}\otimes\idmap_E\big)(\rho).
\end{equation}
In particular $\Crit(\Eprot_M,\rho)=\Crit(\Eprot_F,\sigma)$ for every
\tref{security criterion} in the sense of Definition~\ref{def:criterion}.
\end{theorem}

\begin{proof}
By Lemma~\ref{lem:msmfact}, $\redNmsm_M=\redNmsm_F\circ\Lambda$. Substituting into
\eqref{eq:protocol},
\begin{equation*}
  \Grnd^{M}_j
  =\bar{\Grnd}_j\circ\big(\idmap\otimes\redNmsm_F\circ\Lambda\big)
  =\bar{\Grnd}_j\circ\big(\idmap\otimes\redNmsm_F\big)\circ\big(\idmap\otimes\Lambda_j\big)
  =\Grnd^{F}_j\circ\big(\idmap\otimes\Lambda_j\big),
\end{equation*}
where $\Lambda_j$ denotes $\Lambda$ acting on $\Bsys_j$ and the identity
elsewhere. For $i\ne j$ the maps $\Grnd^{F}_i$ and $\Lambda_j$ act on disjoint
tensor factors --- $\Grnd^{F}_i$ touches $A_i$, $\Qsys_i$ and classical
registers, while $\Lambda_j$ touches only $\Bsys_j$ --- so they commute.
Composing the $n$ rounds therefore moves every $\Lambda_j$ to the front, and
the $\Lambda_j$ assemble into $\Lambda^{\otimes n}$ because they act on
disjoint systems. Applying $\Kfin$, which touches no $\Bsys_j$, gives
\eqref{eq:factorization}. Equation~\eqref{eq:sameoutput} is
\eqref{eq:factorization} tensored with $\idmap_E$, and the equality of the
criteria is \eqref{eq:criterion} evaluated on the two equal states.
\end{proof}
